\section{总结与延伸}

\subsection{讲者的核心总结}

本节课覆盖了两种不同于 RNN 和 Transformer 的文本建模方法：

\begin{enumerate}
    \item \textbf{卷积神经网络}：通过滑动窗口提取局部 n-gram 特征，配合池化操作产生固定长度的句子表示。从简单的 Kim (2014) 单层 CNN 到 VD-CNN 29 层深度字符级网络，展示了 CNN 在文本分类中的潜力。
    \item \textbf{树递归神经网络}：基于语言学的组合性原则，沿句法树自底向上构建短语和句子的语义表示。RNTN 通过张量层实现了对否定等复杂组合语义的建模。
\end{enumerate}

\subsection{全课知识图谱}

\begin{center}
\begin{tikzpicture}[node distance=0.6cm, every node/.style={draw, rounded corners, minimum width=4cm, minimum height=0.7cm, font=\small}]
\node (cnn) {CNN 用于 NLP};
\node (ops) [below=of cnn] {基本操作：卷积、池化、Stride、Dilation};
\node (kim) [below=of ops] {Kim (2014)：多尺度滤波器 + Max Pooling};
\node (vdcnn) [below=of kim] {VD-CNN：深度字符级 CNN};
\node (tree) [below=1.2cm of vdcnn] {树递归神经网络};
\node (rntn) [below=of tree] {RNTN：张量层 + 组合语义};
\node (sst) [below=of rntn] {Stanford Sentiment Treebank};
\draw[->, thick] (cnn) -- (ops);
\draw[->, thick] (ops) -- (kim);
\draw[->, thick] (kim) -- (vdcnn);
\draw[->, thick] (tree) -- (rntn);
\draw[->, thick] (rntn) -- (sst);
\end{tikzpicture}
\end{center}

\subsection{关键 Takeaways}

\begin{importantbox}{五条核心洞见}
\begin{enumerate}
    \item \textbf{CNN 是文本分类的有力工具}：多种大小的滤波器 + Max Pooling 形成了一个简洁高效的文本分类流水线
    \item \textbf{深度和原始信号输入有价值}：VD-CNN 证明了字符级深度 CNN 可以匹敌词级别模型
    \item \textbf{语言的组合语义需要被建模}：关键词匹配无法处理否定、反讽等复杂语言现象
    \item \textbf{结构化模型有其独特优势}：树递归网络在否定建模等方面优于 Transformer，这提示我们结构化的语言学知识仍有价值
    \item \textbf{灵活性与结构之间的权衡}：当前 NLP 的发展趋势是更灵活、更少约束的模型，但这不意味着结构化方法已经过时
\end{enumerate}
\end{importantbox}

\subsection{拓展阅读}

\begin{itemize}
    \item Kim, Y. (2014). Convolutional Neural Networks for Sentence Classification. \textit{EMNLP}. \url{https://arxiv.org/abs/1408.5882}
    \item Conneau, A. et al. (2017). Very Deep Convolutional Networks for Text Classification. \textit{EACL}. \url{https://arxiv.org/abs/1606.01781}
    \item Socher, R. et al. (2013). Recursive Deep Models for Semantic Compositionality Over a Sentiment Treebank. \textit{EMNLP}. \url{https://nlp.stanford.edu/~socherr/EMNLP2013_RNTN.pdf}
    \item Zhang, Y. \& Wallace, B. (2015). A Sensitivity Analysis of (and Practitioners' Guide to) Convolutional Neural Networks for Sentence Classification. \url{https://arxiv.org/abs/1510.03820}
    \item Tai, K.S. et al. (2015). Improved Semantic Representations From Tree-Structured Long Short-Term Memory Networks. \textit{ACL}. \url{https://arxiv.org/abs/1503.00075}
\end{itemize}

\end{document}
